% =============================================================
% 00_howto.tex — 本書の使い方
% =============================================================
\chapter*{この本の使い方}
\addcontentsline{toc}{chapter}{この本の使い方}

\begin{teigi}[ITパスポート試験とは]
ITパスポート試験は、ITに関する基礎的な知識を証明する\textbf{国家試験}です。

\begin{itemize}
  \item \textbf{試験時間}：120分
  \item \textbf{問題数}：100問（全問必須・四肢択一）
  \item \textbf{合格基準}：総合600点以上 ＋ 各分野300点以上
  \item \textbf{出題分野}：ストラテジ系（問1〜35）、マネジメント系（問36〜55）、テクノロジ系（問56〜100）
\end{itemize}
\end{teigi}

\section*{本書の構成}

本書は\textbf{4つのパート}で構成されています。

\begin{hikaku}[4パート構成]
\begin{itemize}
  \item \textbf{Part I}（第1〜2章）：ストラテジ系——企業活動・経営戦略・会計・法務
  \item \textbf{Part II}（第3〜4章）：マネジメント系——開発・PM・サービス・監査
  \item \textbf{Part III}（第5〜7章）：テクノロジ系——基礎理論・NW・DB・セキュリティ・AI
  \item \textbf{Part IV}（第8〜9章）：総仕上げ——横断整理・模擬試験3セット
\end{itemize}
\end{hikaku}

\section*{各章の読み方}

各章は次の流れで構成されています。

\begin{enumerate}
  \item \textbf{この章で身につく力}——学習目標と目安時間
  \item \textbf{要点}——見抜き方のルール
  \item \textbf{例題}——ステップ付きの解法（指針→解答→見抜き方）
  \item \textbf{過去問ドリル}——R6/R7公開問題の演習
  \item \textbf{タイムアタック}——制限時間つきの実戦演習
  \item \textbf{よくあるミス}——ひっかけパターンの対策
  \item \textbf{到達確認チェックリスト}——自己診断
\end{enumerate}

\section*{難易度の表示}

問題には次の難易度が付いています。

\begin{itemize}
  \item \diffbadge{基}\quad 基礎——1つの知識で直接解ける問題
  \item \diffbadge{標}\quad 標準——正しい知識の選択が必要な問題
  \item \diffbadge{強}\quad 強化——複数知識の組合せや、ひっかけを含む問題
  \item \diffbadge{発}\quad 発展——計算や読解力が必要な問題
\end{itemize}

まずは\diffbadge{基}と\diffbadge{標}を確実にしてから、\diffbadge{強}に進みましょう。

\section*{出典について}

\begin{honshiken}[出典表記]
本書に収録する過去問は、独立行政法人情報処理推進機構（IPA）が公表している公開問題に基づきます。
問題文は公開問題の記述を尊重し、学習目的の範囲で引用・参照しています。
各問の冒頭に出典年度を\sourcebadge{R7}のようにバッジで示します。
\end{honshiken}
